Financial literacy remains a significant challenge in India, where large segments of the population lack the knowledge needed to navigate complex financial products such as mutual funds, systematic investment plans and equity instruments. The United Nations Sustainable Development Goals (SDGs) frame financial capability as an integral part of quality education, decent work, reduced inequalities and responsible consumption. This paper presents WealthGenie, an SDG-aligned financial literacy and explainable investment decision support platform designed for Indian beginners. The platform integrates five core modules: a structured Learn curriculum, a Financial Assessment engine, a risk-aware mutual fund recommendation engine, a paper trading simulator, and a Financial Digital Twin for scenario comparison. The recommendation engine is deterministic and explainable, combining verified AMFI and SEBI dataset features with a transparent weighted scoring scheme that evaluates risk alignment, expense ratio, volatility, maximum drawdown, return consistency and data completeness. System evaluation demonstrates implementation consistency, deterministic behaviour and explainability coverage across synthetic profiles. The platform is positioned as an educational tool aligned with SDG 4, SDG 8, SDG 10 and SDG 12, supporting informed financial decision-making without claiming to measure or achieve SDG targets directly.